% ═══════════════════════════════════════════════════════════════
% FORSCHERKARTE: Elektromotor-Experiment (Kl. 9)
% A5 einseitig, zum Laminieren – NUR CHECKLISTEN
% ═══════════════════════════════════════════════════════════════

\documentclass[11pt, a5paper]{article}

\usepackage[utf8]{inputenc}
\usepackage[T1]{fontenc}
\usepackage[ngerman]{babel}

% --- SCHRIFTEN ---
\usepackage{helvet}
\renewcommand{\familydefault}{\sfdefault}

\usepackage{microtype}
\usepackage[margin=1.2cm, top=1.2cm, bottom=1.2cm]{geometry}
\usepackage{tabularx}
\usepackage{array}
\usepackage{amssymb}

\usepackage{xcolor}
\usepackage{tcolorbox}
\tcbuselibrary{skins}

% ═══════════════════════════════════════════════════════════════
% FARBEN
% ═══════════════════════════════════════════════════════════════

\definecolor{physik}{HTML}{2c5aa0}
\definecolor{mittelgrau}{HTML}{888888}
\definecolor{warnung}{HTML}{E65100}
\definecolor{erfolg}{HTML}{2E7D32}

\colorlet{fachfarbe}{physik}

% ═══════════════════════════════════════════════════════════════
% BOXEN – NUR RAHMEN, KEINE FARBFÜLLUNG
% ═══════════════════════════════════════════════════════════════

\newtcolorbox{materialbox}[1][]{
    colback=white,
    colframe=fachfarbe,
    coltitle=black,
    colbacktitle=white,
    boxrule=1.5pt,
    arc=3pt,
    left=6pt, right=6pt, top=4pt, bottom=4pt,
    before skip=5pt, after skip=5pt,
    fonttitle=\bfseries,
    title=#1
}

\newtcolorbox{checklistbox}[1][]{
    colback=white,
    colframe=erfolg,
    coltitle=black,
    colbacktitle=white,
    boxrule=1pt,
    arc=0pt,
    left=6pt, right=6pt, top=4pt, bottom=4pt,
    before skip=5pt, after skip=5pt,
    fonttitle=\bfseries,
    title=#1
}

\newtcolorbox{warnbox}[1][]{
    colback=white,
    colframe=warnung,
    coltitle=black,
    colbacktitle=white,
    boxrule=1.5pt,
    arc=2pt,
    left=6pt, right=6pt, top=4pt, bottom=4pt,
    before skip=5pt, after skip=5pt,
    fonttitle=\bfseries,
    title=#1
}

\newtcolorbox{tippbox}[1][]{
    colback=white,
    colframe=fachfarbe,
    coltitle=black,
    colbacktitle=white,
    boxrule=1pt,
    arc=2pt,
    left=6pt, right=6pt, top=4pt, bottom=4pt,
    before skip=5pt, after skip=5pt,
    fonttitle=\bfseries,
    title=#1
}

% ═══════════════════════════════════════════════════════════════
% BEFEHLE
% ═══════════════════════════════════════════════════════════════

\setlength{\parindent}{0pt}
\setlength{\parskip}{2pt}

\pagestyle{empty}

% ═══════════════════════════════════════════════════════════════
% DOKUMENT
% ═══════════════════════════════════════════════════════════════

\begin{document}

\begin{center}
    {\Large\bfseries\textcolor{fachfarbe}{EXPERIMENTIERKARTE}}\\[0.3em]
    {\normalsize Elektromotor}
\end{center}

\vspace{0.3em}

\begin{materialbox}[Material]
Motor \textbullet{} Netzgerät \textbullet{} Kabel \textbullet{} Stabmagnet \textbullet{} Elektromagnet
\end{materialbox}

\begin{checklistbox}[AUFBAU]
\begin{tabular}{@{}cl@{}}
$\square$ & Netzgerät \textbf{AUS} (0\,V) \\[0.3em]
$\square$ & Motor mittig auf dem Tisch platzieren \\[0.3em]
$\square$ & Kabel an Motor und Netzgerät anschließen \\[0.3em]
$\square$ & Stabmagnet bereithalten \\[0.3em]
$\square$ & \textbf{Erst jetzt:} Spannung langsam hochdrehen \\
\end{tabular}
\end{checklistbox}

\begin{checklistbox}[ABBAU]
\begin{tabular}{@{}cl@{}}
$\square$ & Netzgerät auf \textbf{0\,V} drehen \\[0.3em]
$\square$ & Netzgerät \textbf{ausschalten} \\[0.3em]
$\square$ & Kabel lösen und ordentlich zusammenlegen \\[0.3em]
$\square$ & Magnete in die Box zurücklegen \\[0.3em]
$\square$ & Motor in die Box zurücklegen \\
\end{tabular}
\end{checklistbox}

\begin{warnbox}[SICHERHEIT]
\centering
\begin{tabular}{@{}cl@{}}
$\bullet$ & Maximal \textbf{12\,V} verwenden \\[0.3em]
$\bullet$ & Kabel \textbf{nicht} unter Spannung umstecken \\[0.3em]
$\bullet$ & Motor kann \textbf{heiß} werden -- kurz laufen lassen \\
\end{tabular}
\end{warnbox}

\begin{tippbox}[TIPPS]
\begin{tabular}{@{}p{3.5cm}@{}l@{}}
Nichts passiert? & $\rightarrow$ Kontakte prüfen, Magnet näher \\[0.3em]
Motor brummt nur? & $\rightarrow$ Leicht anschieben, Spannung hoch \\[0.3em]
Dreht rückwärts? & $\rightarrow$ Kabel oder Magnet umpolen \\
\end{tabular}
\end{tippbox}

\vfill

\begin{center}
\scriptsize\textcolor{mittelgrau}{Physik Klasse 9 | Gleichstrommotor}
\end{center}

\end{document}
